\section{An amortized cost model for complete smoothing-criterion workflows}
\label{sec:amortized}
We now distinguish a single mathematical evaluation from a
\emph{workload}. The former concerns the arithmetic complexity of
solving a PLS problem at one fixed penalty; the latter concerns
candidate searches over multiple independent series and multiple
selection criteria. The distinction is crucial because the
smoother matrix and the smoothness-to-penalty map are independent
of the observed series for fixed $(n,d)$.

\subsection{Workload definition and numerical admissibility}
Fix sample length $n$, difference order $1\le d<n$, and write
$r=n-d$. Let $\mathcal{S}=\{s_1,\ldots,s_M\}\subset(0,1)$
be a specified set of $M$ \emph{distinct} interior normalized
smoothness levels, and let $y^{(1)},\ldots,y^{(R)}\in\R^n$
be $R$ observed series. A set $\mathcal C$ of statistical
criteria is evaluated on identical penalties
$\lambda_j=\Lambda_{n,d}(s_j)$.
The criterion label $C$ is allowed to be a singleton
(e.g., CV or GCV) or a set of criteria, so
$c=|\mathcal C|$ need not equal one. Define
\begin{equation}\label{eq:theory_cost_function}
T_{\mathcal A}(n,d,M,R,\mathcal C;\varepsilon,\theta)
=T_{\rm setup}+T_{\rm map}+T_{\rm factors}
+T_{\rm series}+T_{\rm scores}+T_{\rm overhead}.
\end{equation}
Here $\mathcal A$ denotes the complete numerical pipeline,
$\varepsilon$ its requested accuracy tolerances, and $\theta$
the implementation and machine constants. A value of
$T_{\mathcal A}$ is an \emph{execution time}; an asymptotic bound
on its arithmetic work is not an observed number of
milliseconds. We use a single-threaded exact-arithmetic
cost model for proofs, with floating-point admissibility
checked separately.

A pipeline is \emph{admissible} on a prescribed workload only
if every required score is defined, its numerical checks pass,
and its working memory fits the available budget. A small
runtime on a failed or truncated search does not define a
valid winner. This constraint makes the theoretical
frontier a conditional statement rather than a universal
solver ranking.

\subsection{Proposition: reuse of a smoothness grid}
\noindent\textbf{Proposition 1 (data-independent penalty grid).}
For fixed $(n,d)$, the inverse image
$\{\Lambda_{n,d}(s_j)\}_{j=1}^M$ is the same for every
observed series, every fixed numerical difference mean,
and every criterion in $\mathcal C$.

\noindent\emph{Proof.}
Equation~\eqref{eq:S} depends only on the nonzero
eigenvalues of $D_d^\top D_d$. For each finite
$\lambda\ge0$, its derivative is
\[
S_d'(\lambda)=\frac1r\sum_{j=1}^r
\frac{\delta_j}{(1+\lambda\delta_j)^2}>0.
\]
Thus each $s_j\in(0,1)$ has a unique inverse
$\lambda_j$ that is independent of $y$, $\mu$ and
the criterion score. Calculating the inverse twice
for the same $(n,d,s_j)$ cannot change the
mathematical candidate. \hfill$\square$

The proposition does \emph{not} imply that an independently
spaced grid in $\log\lambda$ has the same candidates.
It also does not imply that continuous local optimizers
in $s$ and $\log\lambda$ traverse the same numerical path.

\subsection{Proposition: complete criteria from banded information}
\noindent\textbf{Proposition 2 (banded sufficiency).}
At fixed finite $\lambda>0$ and fixed numeric $\mu$,
an SPD banded Cholesky factorization of
$A_\lambda=I+\lambda D_d^\top D_d$, together with
the selected inverse on its sparsity band, suffices
to evaluate the fitted curve, conditional
CV/GCV/AICc/residual BIC, and a profiled Gaussian
reduced-coordinate marginal-likelihood score,
without forming a dense $n\times n$ inverse.

\noindent\emph{Proof.}
The direct normal equations are
\begin{equation}\label{eq:theory_rhs}
A_\lambda\widehat\tau_\lambda
=y+\lambda\mu D_d^\top\mathbf1.
\end{equation}
One banded triangular solve gives $\widehat\tau$.
A backward selected-inverse recursion
(Section~\ref{sec:structured}) gives
$h_i=(A_\lambda^{-1})_{ii}$, and hence
$\nu=\sum_i h_i$. Residuals $e=y-\widehat\tau$
and $\mathrm{RSS}=e^\top e$ then determine
CV, GCV, AICc and residual BIC through
Eqs.~\eqref{eq:cv_gcv}--\eqref{eq:information}.

For the Gaussian reduced-coordinate likelihood,
choose $p_\mu$ with $D_dp_\mu=\mu\mathbf1$, and put
$z=y-p_\mu$. Let $w_j=\lambda\delta_j/(1+\lambda\delta_j)$.
Then
\begin{equation}\label{eq:theory_profile_quad}
\sum_{j=1}^r w_j(v_j^\top z)^2
=z^\top(I-H_\lambda)z
=\|y-\widehat\tau_\lambda\|^2
+\lambda\|D_d\widehat\tau_\lambda-\mu\mathbf1\|^2
=:J_\lambda.
\end{equation}
The final equality follows by substituting the
normal-equation minimizer in the strictly convex
penalized objective. Thus $J_\lambda$ is available
from the fitted curve and a banded difference pass;
there is no need to construct the potentially
ill-scaled polynomial $p_\mu$.

For $\lambda>0$, the marginal variance multipliers are
$c_j=1+1/(\lambda\delta_j)$. Since the $r$ nonzero
eigenvalues of $Q_d$ are the eigenvalues of
$D_dD_d^\top$, the determinant identity is
\begin{equation}\label{eq:theory_logdet}
\begin{aligned}
\sum_{j=1}^r\log c_j
&=\sum_{j=1}^r\log(1+\lambda\delta_j)
-r\log\lambda-\sum_{j=1}^r\log\delta_j\\
&=\log\det A_\lambda-r\log\lambda
-\log\det(D_dD_d^\top).
\end{aligned}
\end{equation}
The first log determinant is the sum of twice the
logarithms of the diagonal Cholesky entries
of $A_\lambda$. The last determinant is independent
of $\lambda$, $y$ and $\mu$ and can be prepared
once using a banded factorization of
$D_dD_d^\top$. Profiling the observational variance
gives $\widehat\sigma^2_\varepsilon=J_\lambda/r$
and a negative twice-log-likelihood, up to constants,
\begin{equation}\label{eq:theory_mle}
\mathcal L(\lambda)
=r\log(J_\lambda/r)
+\log\det A_\lambda-r\log\lambda
-\log\det(D_dD_d^\top).
\end{equation}
Consequently all the stated quantities can be
evaluated from banded information alone.
\hfill$\square$

The likelihood statement assumes the latent-difference
Gaussian variance model of Eq.~\eqref{eq:gauss}.
It is not ordinary residual least squares and is
not joint likelihood optimization over $\mu$.
For $\mu=\operatorname{mean}(D_dy)$ computed from the
same series, the fixed-$\mu$ leverage formula is
\emph{conditional}, not an exact leave-one-out
refit of the plug-in estimator. At $\lambda=0$
some selection formulas are undefined and the
Gaussian prior limit is improper; $s=0$ and $s=1$
must be treated by their own endpoint conventions.

\subsection{Theorem: arithmetic and memory bounds}
We assume a conventional full retained SVD of $D_d$,
an SPD banded factor and Takahashi-type selected
inversion, and fixed $d$-banded operations.
Set $\chi_{\mathrm{CV}}=1$ when ordinary
conditional CV belongs to $\mathcal C$ and zero
otherwise. Let $I$ be a uniform upper bound on
the number of function evaluations used to
invert each interior $s_j$ to the required
scalar tolerance.

\noindent\textbf{Theorem 1 (amortized work).}
With reusable full spectral preprocessing and a
single shared smoothness-to-penalty map, the
arithmetic work for $R$ series and $M$ candidates
is bounded by
\begin{equation}\label{eq:theory_spectral_bound}
W_{\rm SVD}
=O\bigl(nr^2+MIr+Rnr+MRcr
+\chi_{\rm CV}MRnr\bigr).
\end{equation}
Its retained dense-basis memory is $O(nr)$, in
addition to $O(M+Rn)$ if all the observed
series are held simultaneously.

Using banded Cholesky and exact selected-inverse
diagnostics, and evaluating $S_d(\lambda)$
by banded trace calculations during root finding,
the arithmetic work is bounded by
\begin{equation}\label{eq:theory_banded_bound}
W_{\rm band}
=O\bigl(nd^2+MI\,nd^2+Mnd^2
+MRnd+MRcn\bigr).
\end{equation}
When the penalty grid has already been mapped,
the $MI\,nd^2$ term is absent. A streamed
implementation retains $O(nd+M+Rn)$ storage;
retaining \emph{all} $M$ banded factors instead
requires $O(Mnd)$ storage.

\noindent\emph{Proof.}
The reduced SVD has cost $O(nr^2)$ and
stores $O(nr)$ entries. Each spectral
smoothness evaluation is a sum of $r$
scalars, so $M$ inversions of at most $I$
evaluations cost $O(MIr)$. One projection
per observed series requires $O(Rnr)$.
Once projected, any of the $c$ scalar
scores at each $(j,k)$ is $O(r)$.
If CV is included, exact leverage and
full residual reconstruction are bounded
by $O(nr)$ per candidate and series,
giving the additional CV term.
This is a conservative bound; leverage
vectors can be shared among the $R$
series at each candidate.

For the banded method, preparation of the
banded difference-penalty matrix and the
constant Gaussian determinant costs
$O(nd^2)$. Every trace-based root-solver
evaluation factors an SPD matrix and
computes the needed selected inverse in
$O(nd^2)$ operations. After mapping,
factorization and selected inversion are
done once per candidate in $O(Mnd^2)$,
reused across all criteria and all
observed series. For each candidate
and series, two triangular solves,
difference passes and quadratic summaries
use $O(nd)$ arithmetic, followed by
at most $O(cn)$ score formation.
Processing one candidate at a time
requires one banded factor and selected
inverse; keeping all observations
uses $O(Rn)$ additional storage.
Caching all factors produces the
stated $O(Mnd)$ alternative.
\hfill$\square$

These are upper bounds for explicit, fully specified
implementation families. They do not prove
a lower bound for every other possible algorithm;
fast transforms for selected difference orders
or alternative numerical approximations can
produce different costs. Furthermore, a large
root-iteration count $I$ can erase the benefit of
banded factorization when the user insists on
inverting an extremely fine $s$ grid.

\subsection{Corollary: order-one difference spectra}
\noindent\textbf{Corollary 1.}
For the first-difference matrix with the
ordinary finite-interval boundary convention,
\begin{equation}\label{eq:theory_dct}
\delta_j=2-2\cos(\pi j/n),\qquad j=1,\ldots,n-1.
\end{equation}
Thus the entire positive penalty spectrum
is available in $O(n)$ arithmetic and
storage without a dense SVD, and the
$M$-point smoothness grid can be inverted
in $O(MIn)$ arithmetic through direct
spectral sums.

\noindent\emph{Proof.}
$D_1^\top D_1$ is the path-graph
Laplacian with endpoint diagonal values
one, interior diagonal values two, and
adjacent off-diagonals $-1$. Its cosine
eigenvectors have eigenvalues
$2-2\cos(\pi j/n)$, $j=0,\ldots,n-1$.
The zero eigenvalue corresponds to the
constant vector. Summing the remaining
$n-1$ eigenvalues in Eq.~\eqref{eq:S}
takes $O(n)$ arithmetic per root evaluation.
\hfill$\square$

This special case is an established discrete-cosine
identity, not a claimed new spectral algorithm.
For general $d>1$, endpoint effects mean that
the same elementary eigenvalue expression
cannot simply be substituted.

\subsection{When can a crossover exist?}
\label{sec:crossover}
A method crossover depends on the machine,
statistical workload, numerical tolerance and
actual pipeline implementation. To expose
this dependence, write a local empirical
calibration model for a \emph{fixed}
$(n,d,R,\mathcal C,\varepsilon,\theta)$:
\begin{equation}\label{eq:theory_affine_time}
T_{\mathcal A}(M)=a_{\mathcal A}
+b_{\mathcal A}M+\rho_{\mathcal A}(M),
\end{equation}
where $a_{\mathcal A}\ge0$ is fixed preparation,
$b_{\mathcal A}\ge0$ is per-candidate effective
work at that workload, and
$\rho_{\mathcal A}(M)$ records cache changes,
allocation, nonlinear solver iterations
and contention effects. This is a
\emph{calibration ansatz}, not an
exact equality implied by big-$O$.

\noindent\textbf{Proposition 3 (conditional crossover).}
If $\rho_{\rm SVD}=\rho_{\rm band}=0$
over a declared workload interval,
$a_{\rm SVD}>a_{\rm band}$ and
$b_{\rm band}>b_{\rm SVD}$, then banded
Cholesky is faster for
$M<M_\star$, spectral SVD is faster for
$M>M_\star$, and
\begin{equation}\label{eq:theory_threshold}
M_\star=
\frac{a_{\rm SVD}-a_{\rm band}}
{b_{\rm band}-b_{\rm SVD}}.
\end{equation}
If the sign conditions do not hold, no
positive crossover of that direction
follows from the affine model.

\noindent\emph{Proof.}
Subtract the two affine time functions:
\[
T_{\rm SVD}(M)-T_{\rm band}(M)
=(a_{\rm SVD}-a_{\rm band})
-M(b_{\rm band}-b_{\rm SVD}).
\]
Its unique zero is $M_\star>0$ under
the stated inequalities, and its
sign changes once at this point.
\hfill$\square$

The proposition does \emph{not} supply a
universal observed threshold: constants
must be calibrated on comparable,
numerically admissible workloads.
When ordinary CV is required, the
$O(nr)$ spectral candidate work may
change the signs. When multiple series
share an already-mapped $s$ grid and
reuse banded factors, the effective
$b_{\rm band}$ can decrease. Consequently,
a universal difference-order-only
dispatcher is not justified by these
complexity bounds.

\subsection{Finite-memory frontier and candidate streaming}
For an eight-byte floating-point type,
one dense $n\times n$ matrix occupies
$8n^2$ bytes before copies or solver
workspaces. At $n=131\,072$ this is
$128$ GiB, so an ordinary dense inverse
can fail the memory budget before any
statistical criterion is evaluated.
Banded factorizations instead require
$O(nd)$ storage for fixed $d$.
An implementation may stream penalty
candidates one at a time, updating the
best valid criterion score per series
and criterion, without retaining $M$
full smoothing matrices.

The total-time functional
\eqref{eq:theory_cost_function} should
therefore be minimized over the
\emph{admissible set} of algorithms:
\begin{equation}\label{eq:theory_frontier}
\mathcal A^\star
\in\arg\min_{\mathcal A\in\mathfrak A_{\rm valid}}
T_{\mathcal A}(n,d,M,R,\mathcal C;\varepsilon,\theta).
\end{equation}
Here $\mathfrak A_{\rm valid}$ is
workload-specific: statistical
undefinedness, forward-error limits,
numerical nonconvergence and the memory
budget can remove a method from the
comparison. The formula is a
decision framework and not a claim
that its minimizer has been
empirically identified across the
full design space.
